\problemname{Train Stations}

Zohan and Jimón are on their way to a programming training camp. The epic training camp takes place in the city of Petrozavodsk, and they have decided to travel by train.

During the journey, Jimón is for some reason counting the number of people who get on and off at each station the train stops at. He writes these numbers
down in his notebook -- one station's data per page.

When Jimón is about to get off the train he stumbles and his notebook is torn to pieces -- all he has left is a pile of notes in a jumbled order. Zohan now challenges
Jimón to reconstruct the order in which the stations appeared, given the numbers written on the pages lying on the ground. Can you help him, or can you prove that
Jimón must have miscounted?

\section*{Input}
The first line contains an integer $N$, the number of pages in the notebook.

After that follow $N$ lines (one per page), each with two non-negative
integers: the number of people who boarded at the station, and the number who
got off.

A person never gets off at the same station where they boarded. It is guaranteed that
the total number of boardings and the total number of
alightings are the same, and that this number is at most $10^9$. The train is always empty
when Jimón starts counting and the train is always empty when he has finished counting. For
simplicity we assume that Jimón does not count himself or Zohan -- and
it is guaranteed that they board first and alight last (so they never miss
counting anyone).

\section*{Output}
If it is possible for the stations to be ordered so that there are never more
people alighting than are on the train, first output a line
``\texttt{JA}'', and then a line with a possible order, where each number $1$ to
$N$ appears exactly once.
Otherwise, i.e. if Jimón has made a mistake, output ``\texttt{NEJ}''.

\section*{Explanation of Samples}
In the first example, a total of $12$ people ride the train. At the first station
4 people board and nobody alights, then the train continues to
the next station. There an additional $8$ board and $2$ alight, so the train now has $10$
people on board. At the last station the remaining $10$ alight. This is described by the order
\texttt{2 3 1}.

In the second example we know that the first station must be either the one where $2$ board or where $3$
board. Thinking further, we realize that these two must be first if there is to be any valid solution,
which eventually leads us to the conclusion that it will not be possible to create any solution without too many
people alighting. Note that \emph{2 3 4 1} is not a valid solution since at station four we would need a person
to alight at the same station where they boarded -- which is not allowed.

\section*{Scoring}
Your solution will be tested on a set of test groups. To earn points for a group, you must pass all test cases in that group.

\begin{tabular}{| l | l | l | l |}
\hline
Group & Points & Constraints & Other\\ \hline
1     & 20         & $ 2 \le N \le 8 $ & \\ \hline
2     & 20         & $ 2 \le N \le 1\,000 $, there will only be one station where more people alight than board & \\ \hline
3     & 30         & $ 2 \le N \le 1\,000 $ & \\ \hline
4     & 30         & $ 2 \le N \le 60\,000 $ & \\ \hline
\end{tabular}
